\subsection{The Parametric Equation of a Line}

The equations used so far describe a line by a condition that its coordinates
must satisfy. A parametric equation describes the same line differently: we
start from one point and move along a fixed direction.

Let
\[
P_0=(x_0,y_0)
\]
be a point of the line, and let
\[
(u,v)\neq(0,0)
\]
be a direction pair. Multiplying the direction pair by a real number $t$
changes only the distance and orientation of the displacement, not its
direction. Starting at $P_0$ therefore gives the points
\[
(x,y)=(x_0,y_0)+t(u,v),
\qquad t\in\mathbb{R}.
\]

\begin{definitionbox}
The line through $P_0=(x_0,y_0)$ with direction pair $(u,v)\neq(0,0)$ is
parametrised by
\[
L=\{(x_0+tu,\,y_0+tv):t\in\mathbb{R}\}.
\]
Equivalently, its coordinates satisfy
\[
\begin{cases}
x=x_0+tu,\\
y=y_0+tv,
\end{cases}
\qquad t\in\mathbb{R}.
\]
\end{definitionbox}

When $t=0$, the parametrisation gives the initial point $P_0$. Positive and
negative values of $t$ move from $P_0$ in the two opposite directions along the
line. Since $t$ ranges over all real numbers, every point of the line is
obtained.

\begin{center}
\begin{tikzpicture}[scale=0.9]
    \draw[axis] (-1.0,0) -- (6.2,0) node[right] {$x$};
    \draw[axis] (0,-1.0) -- (0,4.8) node[above] {$y$};

    \coordinate (P) at (1.2,1.1);
    \coordinate (Q) at (4.8,3.5);
    \coordinate (D) at (3.0,2.3);

    \draw[subset] (0.3,0.5) -- (5.6,4.03) node[above right] {$L$};
    \fill[geometricSubsetColor] (P) circle (2.2pt)
        node[below right] {$P_0=(x_0,y_0)$};
    \fill[geometricSubsetColor] (Q) circle (2.2pt)
        node[above left] {$P_0+t(u,v)$};
    \draw[->,helper] (P) -- (D)
        node[midway,above left] {$(u,v)$};
\end{tikzpicture}
\end{center}

If $u\neq0$, eliminating the parameter gives
\[
t=\frac{x-x_0}{u}
\]
and hence
\[
y-y_0=\frac{v}{u}(x-x_0).
\]
Thus the slope is
\[
m=\frac{v}{u}.
\]
If $u=0$, then $v\neq0$ and the parametrisation becomes
\[
x=x_0,
\qquad
y=y_0+tv,
\]
which describes a vertical line. The parametric form therefore treats
non-vertical and vertical lines in one uniform way.

For two distinct points
\[
P=(x_1,y_1),
\qquad
Q=(x_2,y_2),
\]
the difference
\[
Q-P=(x_2-x_1,\,y_2-y_1)
\]
is a direction pair for their line. Consequently,
\[
L=\{P+t(Q-P):t\in\mathbb{R}\},
\]
or in coordinates,
\[
\begin{cases}
x=x_1+t(x_2-x_1),\\
y=y_1+t(y_2-y_1),
\end{cases}
\qquad t\in\mathbb{R}.
\]

\begin{examplebox}
The line through $P=(1,2)$ with direction pair $(3,6)$ has the parametrisation
\[
\begin{cases}
x=1+3t,\\
y=2+6t,
\end{cases}
\qquad t\in\mathbb{R}.
\]
Because $3\neq0$, we have
\[
t=\frac{x-1}{3}.
\]
Substitution into the second equation gives
\[
y=2+6\frac{x-1}{3}=2x.
\]
Thus the parametric and Cartesian equations describe the same line.
\end{examplebox}

\begin{remarkbox}
A parametrisation of a line is not unique. Replacing $(u,v)$ by any non-zero
multiple $\lambda(u,v)$ changes the parameter but not the set of points on the
line.
\end{remarkbox}
